% Ad Familiares 16.10 --- headnote
% ad-familiares-16-10, 17 April 53 BC, Cumae

\textit{Cicero to Tiro, written from the villa at Cumae on the
fourteenth day before the Kalends of May --- 17 April 53 BC.
Tiro's fever has broken; he is well enough that Cicero now
worries about the relapse a difficult journey could bring on
rather than about the illness itself. Acastus, the slave Cicero
had urged Tiro to keep on as a nurse (see 16.14), has at last
brought a letter back. The patron's calculation in the first
section is touchingly precise: two days on the road, five for
the return trip, all measured against the date Cicero wants to
be at Formiae --- the third day before the Kalends.}

\textit{The second section opens with one of the warmest jokes
in the correspondence: Cicero's \textit{litterulae meae sive
nostrae}, his ``little pages, or rather ours,'' have grown
listless without their secretary, and only Acastus's delivery
has raised their eyes again. Pompey is in the house as he
writes, asking after the work, and Cicero tells him plainly
that without Tiro everything of his is mute. The closing
promise --- \textit{nostra ad diem dictam fient} --- is the
manumission already pencilled in for the day Tiro arrives; the
clinching tag, that he has long since taught Tiro the
\textit{etymon} of \textit{fides}, is the small Greek word
Cicero leaves in transliterated script and Latin grammar at
once, a private joke between two men used to working in both
languages.}
